\documentclass[11pt]{beamer}
\usepackage[utf8]{inputenc}
\usepackage[T1]{fontenc}
\usepackage{booktabs} 

%%% Customize Theme %%%%%%%%%%%%%%%%%%%%%%
\usetheme{Madrid} 

\definecolor{myRed}{RGB}{134,31,65}
\definecolor{myOrange}{RGB}{229, 117, 31}

\setbeamercolor*{structure}{bg=myRed!20,fg=myRed!90}
\setbeamercolor*{palette primary}{use=structure,fg=white,bg=structure.fg} 
\setbeamercolor*{palette tertiary}{use=structure,fg=white,bg=myRed}
\setbeamercolor{frametitle}{bg=myRed!85,fg=white} 
\setbeamercolor*{titlelike}{parent=palette primary} 
\setbeamercolor{section in head/foot}{fg=myOrange, bg=white}

%%% Specific Colors %%%
\setbeamercolor{item projected}{bg=myOrange}
\setbeamertemplate{enumerate items}{bg=myOrange}
\setbeamercolor{itemize item}{fg=myOrange}
\setbeamercolor{itemize subitem}{fg=myOrange}
\setbeamercolor{button}{bg=myOrange}
\setbeamercolor{section in toc}{fg=black}
\setbeamercolor{subsection in toc}{fg=black}

%%% Block Colors %%%
\setbeamercolor{block title}{bg=myOrange, fg=white}
\setbeamercolor{block body}{bg=myOrange!20}

%---------------------------------------------------------
%	SELECT FONT THEME & FONTS
%---------------------------------------------------------
\usefonttheme{default} 
\usepackage{palatino} 
\usepackage[default]{opensans} 
\useinnertheme{circles}

%---------------------------------------------------------
%	PRESENTATION INFORMATION
%---------------------------------------------------------
\title[Composants Électroniques]{Les 4 composants électroniques de base}
\subtitle{Résistance, Condensateur, Diode, Transistor}
\author[Walid EL BEKRAOUI]{Walid EL BEKRAOUI}

\date[\today]

\begin{document}

%---------------------------------------------------------
%	TITLE SLIDE
%---------------------------------------------------------
\begin{frame}
    \titlepage
\end{frame}

%---------------------------------------------------------
%	TABLE OF CONTENTS SLIDE
%---------------------------------------------------------
\begin{frame}
    \frametitle{Sommaire}
    \tableofcontents
\end{frame}

%---------------------------------------------------------
%	PRESENTATION BODY SLIDES
%---------------------------------------------------------

\section{Objectifs pédagogiques}
\begin{frame}
    \frametitle{Objectifs pédagogiques}
    À l'issue de cette formation, vous serez capable de :
    \begin{itemize}
        \item Identifier ces composants visuellement et par leur symbole
        \item Expliquer leur rôle dans un circuit
        \item Savoir lire les valeurs (code couleurs, marquages)
        \item Différencier les principaux types
        \item Comprendre leurs applications concrètes
        \item Effectuer un test simple avec un multimètre
    \end{itemize}
\end{frame}

\section{La résistance}
\begin{frame}
    \frametitle{1. La résistance}
    \begin{block}{1.1 Définition}
        La résistance est un composant passif qui limite l'intensité du courant électrique qui circule dans un circuit. Elle transforme une partie de l'énergie électrique en chaleur.
    \end{block}
    
    \textbf{1.2 Loi fondamentale (Loi d'Ohm)}
    $$ U = R \times I $$
    \begin{itemize}
        \item $U$ : Tension en Volts (V)
        \item $R$ : Résistance en Ohms ($\Omega$)
        \item $I$ : Intensité en Ampères (A)
    \end{itemize}
\end{frame}

\begin{frame}
    \frametitle{La résistance : Types et Applications}
    \textbf{1.6 Types de résistances}
    \begin{itemize}
        \item \textbf{Carbone / Métallisée :} Circuits génériques, audio, mesures
        \item \textbf{Potentiomètre :} Résistance variable (volume)
        \item \textbf{Thermistance / LDR :} Varie avec température / lumière
    \end{itemize}
    
    \vspace{0.5cm}
    \textbf{1.7 Applications concrètes}
    \begin{itemize}
        \item Limiter le courant (ex: Protéger une LED)
        \item Diviseur de tension
        \item Pull-up / Pull-down (Fixer un niveau logique)
    \end{itemize}
    
    \begin{alertblock}{Exercice n°1}
        Calculez la résistance nécessaire pour alimenter une LED rouge ($V_F = 1,7V$, $I = 10mA$) sous $5V$. Astuce : $R = \frac{U_{alim} - V_F}{I}$
    \end{alertblock}
\end{frame}

\section{Le condensateur}
\begin{frame}
    \frametitle{2. Le condensateur (Capacité)}
    \begin{block}{2.1 Définition}
        Le condensateur est un composant passif capable de stocker et de restituer de l'énergie électrique sous forme de charges. Il agit comme un réservoir temporaire de tension.
    \end{block}
    
    \textbf{2.2 Grandeur et unité}
    La capacité $C$ se mesure en Farads (F). 
    \begin{itemize}
        \item Microfarad ($\mu F$) : $10^{-6}F$
        \item Nanofarad ($nF$) : $10^{-9}F$
        \item Picofarad ($pF$) : $10^{-12}F$
    \end{itemize}
\end{frame}

\begin{frame}
    \frametitle{Le condensateur : Comportement et Polarité}
    \textbf{2.5 Comportement}
    \begin{itemize}
        \item \textbf{Continu (CC) :} Se comporte comme un circuit ouvert après charge.
        \item \textbf{Alternatif (CA) :} Laisse passer le courant (plus la fréquence est élevée, plus il se laisse traverser).
    \end{itemize}
    
    \vspace{0.3cm}
    \begin{alertblock}{2.7 Polarité à respecter (TRÈS IMPORTANT)}
        Un condensateur électrolytique inversé peut gonfler, fuir ou exploser.
        \begin{itemize}
            \item Borne + : Souvent plus longue ou marquée d'un "+"
            \item Borne -- : Marquée d'une bande avec des "--"
        \end{itemize}
    \end{alertblock}
\end{frame}

\section{La diode}
\begin{frame}
    \frametitle{3. La diode}
    \begin{block}{3.1 Définition}
        La diode est un composant semi-conducteur qui ne laisse passer le courant électrique que dans un seul sens. C'est un « clapet » électronique unidirectionnel (Jonction PN).
    \end{block}

    \textbf{3.3 Polarisation}
    \begin{itemize}
        \item \textbf{Passante :} Anode (+) vers Cathode (-) $\rightarrow$ Conductrice.
        \item \textbf{Bloquante :} Inversée $\rightarrow$ Isolante.
    \end{itemize}

    \textbf{3.5 Tension de seuil typique ($V_F$)}
    \begin{itemize}
        \item Silicium (ex: 1N4148) : $\sim 0,7V$
        \item LED Rouge : $\sim 1,6 - 1,8V$
    \end{itemize}
\end{frame}

\section{Le transistor}
\begin{frame}
    \frametitle{4. Le transistor}
    \begin{block}{4.1 Définition}
        Composant semi-conducteur actif. Son rôle est de commander un courant important à l'aide d'un très faible courant. (Interrupteur commandé ou amplificateur).
    \end{block}
    
    \begin{columns}[t]
        \begin{column}{0.5\textwidth}
            \textbf{4.3 Constitution (BJT)}
            \begin{itemize}
                \item \textbf{B}ase : Commande
                \item \textbf{É}metteur : Sortie
                \item \textbf{C}ollecteur : Courant principal
            \end{itemize}
        \end{column}
        \begin{column}{0.5\textwidth}
            \textbf{4.4 Les 2 types}
            \begin{itemize}
                \item \textbf{NPN} : $I_C \to E$
                \item \textbf{PNP} : $E \to I_C$
            \end{itemize}
        \end{column}
    \end{columns}
\end{frame}

\begin{frame}
    \frametitle{Le transistor : Fonctionnement}
    \textbf{4.5 États de fonctionnement}
    \begin{itemize}
        \item \textbf{Coupure (OFF) :} $I_B = 0 \rightarrow I_C \approx 0$ (Interrupteur ouvert)
        \item \textbf{Saturation (ON) :} $I_B$ suffisant $\rightarrow I_C$ maximal (Interrupteur fermé)
        \item \textbf{Zone active :} $I_C = \beta \times I_B$ (Amplificateur)
    \end{itemize}
    
    \vspace{0.5cm}
    \textbf{4.7 Exemple : Piloter une LED (NPN)}
    \begin{itemize}
        \item Bouton relâché $\rightarrow I_B = 0 \rightarrow$ Transistor coupé $\rightarrow$ LED éteinte
        \item Bouton appuyé $\rightarrow$ Courant dans la base $\rightarrow$ Transistor saturé $\rightarrow$ LED allumée
    \end{itemize}
\end{frame}

\section{Comparaison synthétique}
\begin{frame}
    \frametitle{5. Comparaison synthétique}
    \footnotesize
    \begin{table}[]
        \begin{tabular}{@{}lllll@{}}
        \toprule
        \textbf{Composant} & \textbf{Type} & \textbf{Rôle principal} & \textbf{Unidir.} & \textbf{Stockage} \\ \midrule
        Résistance & Passif & Limiter courant / Diviser & Non & Non \\
        Condensateur & Passif & Stocker / Restituer charge & Non* & Oui \\
        Diode & Semi-cond. & Redresser / Protéger / LED & Oui & Non \\
        Transistor BJT & Actif & Interrupteur / Amplificateur & Dépend & Non \\ \bottomrule
        \end{tabular}
    \end{table}
    \vspace{0.5cm}
    \textit{*Sauf effet polarisé pour les électrolytiques/tantales.}
\end{frame}

%---------------------------------------------------------
%	CLOSING SLIDE
%---------------------------------------------------------
\appendix
\begin{frame}[noframenumbering]
    \begin{center}
        {\LARGE Merci de votre attention !\\Avez-vous des questions ?}
    \end{center}
\end{frame}

\end{document}
